\section*{Code and Data Availability}

The code, the experiment outputs, and a Jupyter notebook that reproduces every table and figure of this paper are available at \url{https://github.com/nna17/wifi-rtt-rss-feature-subsampling} and archived at \textbf{[Zenodo DOI]} (release \textbf{[tag, e.g., v1.0]}). The repository contains the main experiment pipeline, the extension and screening experiments, the scripts that generate all tables and figures from the experiment outputs, the exact package versions used (Python 3.14, scikit-learn 1.9.1), and step-by-step instructions; with the included outputs, the tables and figures can be regenerated in seconds without rerunning the experiments. The notebook additionally checks that every number reported in this paper matches the experiment outputs.

The hybrid RTT--RSS benchmark of Feng et al.~\cite{feng2026robust} is available at \url{https://doi.org/10.5281/zenodo.11558192} and \url{https://github.com/Fx386483710/WiFi-RTT-RSS-dataset}. The EETAC auditorium dataset~\cite{martinescalona2023ftm} is available at \url{https://github.com/imartin-UPC/RTT-Data-Collection}. The datasets are not redistributed with the code; the repository describes where to place them.
